\paragraph{Crossovers are early, when they happen at all.}
For three of the four pairs the asymptotically faster algorithm wins from
inputs that fit comfortably in L1 cache: 1\,500 keys, a $32\times32$ matrix,
a 1\,100-bit integer. The common belief that ``the fancy algorithm only
pays off for huge inputs'' is, for these three, off by several orders of
magnitude. What the belief gets right is the fourth pair: an algorithm whose
advantage lies in an operation that realistic inputs rarely perform, and whose
implementation trades array traversal for pointer chasing, can fail to win at
any size that fits in memory.

\paragraph{The parameter moves the crossover more than the exponent does.}
For Karatsuba the base-case cutoff shifts $n^*$ by a factor of twenty (34
limbs at cutoff 32, 700 at cutoff 8) and halves the asymptotic advantage; for
Strassen it moves $n^*$ between 32 and roughly 1\,000 and changes the speed at
$n=4096$ by $1.6\times$; the shape of the matrix---how many factors of two
its dimension has---moves Strassen's crossing by a factor of ten; and the loop
order of the \emph{baseline} is worth $17\times$ at $n=2048$ where
Strassen's exponent is worth $2.6\times$. In the measured range, engineering decisions that the
asymptotic analysis cannot see dominate the decision it can. The practical
reading is that a recursive algorithm is only as good as its base-case policy,
and that a ``naive'' baseline must be a competent one before any comparison
means anything.

\paragraph{Wide intervals are answers, not failures.}
Several crossings in Table~\ref{tab:crossovers} have intervals spanning a
factor of two or more: Strassen with a too-small cutoff, Karatsuba with a
too-small cutoff, radix sort with too-wide digits. In each the ratio curve approaches 1 slowly, and the
bootstrap faithfully reports that the crossing could be anywhere in a wide
range. For a practitioner this is the useful fact: over that range the two
implementations are within noise of each other and the simpler one should be
preferred for its simplicity; the asymptotic winner only earns its complexity
beyond the top of the interval, and then only by the modest ratio we report
there. Sharp crossings, by contrast, are where the decision is worth making
carefully, and the sorting and Karatsuba (cutoff 32) rows show the grid can
resolve them to $\pm10\%$ and to a few limbs respectively.

\paragraph{Input structure can flip the verdict.}
Radix sort's loss on nearly-sorted and reverse-sorted input and the Fibonacci
heap's parity only on complete graphs are the same phenomenon seen from two
sides. An algorithm
with a worse bound but adaptive behaviour (introsort, the lazy heap) exploits
structure that the input happens to have; an algorithm with a better bound
and oblivious behaviour (radix sort, the Fibonacci heap's decrease-key) does
exactly the same work regardless. Whether that is a feature or a bug depends
on whether the structure is present, and the atlas gives the practitioner a
way to check: the key is the input family, not the size.

\paragraph{Measured exponents are honest in the middle and lie at the edges.}
Over the largest decade of each sweep the fitted exponents sit close to theory
for the compute-bound algorithms (Karatsuba 1.55--1.60, schoolbook
1.97--2.01, Strassen 2.77--2.86) and drift above it for the memory-bound ones (introsort 1.14 on 64-bit keys where $n\log n$ predicts $\approx1.06$
over that decade; the $i$-$j$-$k$ loop 4.3 where theory says 3). A fitted exponent on a log--log plot is a statement about the cache
hierarchy as much as about the algorithm, which is why the cache boundaries
are drawn on every runtime figure.

\paragraph{What the atlas cannot say.}
These are single-threaded scalar implementations on one machine, measured
under a hypervisor without hardware counters. A vectorised radix sort or a
BLAS kernel would move every crossover, and not necessarily in the direction one
expects, since vectorisation may help the simple kernel more than the
recursive one built on it. The harness is
offered as the way to redraw the map rather than the map as a permanent
answer; Section~\ref{sec:repro} gives the one command that does so.

\paragraph{Guidance.}
For a practitioner choosing between these pairs today, on hardware like
ours:
\begin{itemize}\itemsep0pt
  \item Sort with radix sort (11-bit digits) when keys are integers, inputs
  exceed about 1\,500 elements, and the data is not already nearly or
  reverse sorted; use \texttt{std::sort} otherwise.
  \item Use Strassen with a base case of 32--64 for dense matrices above
  about 64 rows when the dimension is a multiple of a high power of two,
  or above about 300 otherwise (or pad once at the top)---and fix the loop
  order of the baseline first.
  \item Use Karatsuba above about 32 limbs with a cutoff of 32--128; the
  libraries already do.
  \item Do not use a Fibonacci heap for Dijkstra on any graph that fits in
  memory. Use a binary heap with lazy deletion unless the graph is a grid-like
  network, where an indexed binary heap is slightly better.
\end{itemize}
